\chapter{Food Safety}

\section*{Definition and Overview}
Food safety is the set of practices that prevent food from becoming contaminated and causing illness. Unsafe food can contain bacteria, viruses, parasites, chemicals, or foreign objects that make people sick, and the risks exist at every step from the shop to the table: storing, preparing, cooking, and serving food. The core rules of food safety are simple and memorable: clean, separate, cook, and chill. Cleaning hands and surfaces, separating raw and cooked foods, cooking to safe temperatures, and chilling food promptly prevent nearly all foodborne illness. Food safety matters most for vulnerable people, including children, older adults, pregnant women, and people with weakened immunity. This chapter explains the rules and how to apply them at home.

\section*{Epidemiology}
Foodborne illness affects around 4 million Australians each year, with an estimated 1 in 5 people affected annually. Most cases occur in the home, often from undercooked meat and poultry, inadequate refrigeration, cross-contamination, and poor hand hygiene. Salmonella, campylobacter, and norovirus are leading causes, and outbreaks are traced to homes, restaurants, and community events. Vulnerable groups account for most hospitalisations and deaths. Population surveys show that most people believe their own food handling is safe, yet common mistakes, such as not cooking poultry thoroughly and washing raw chicken, remain widespread.

\section*{Aetiology and Pathophysiology}
\begin{itemize}
    \item \textbf{Bacteria:} salmonella and campylobacter (poultry, eggs, meat), listeria (deli meats, soft cheeses), E.\ coli (minced beef, produce), and staphylococcus (prepared foods)
    \item \textbf{Viruses:} norovirus from infected food handlers, and hepatitis A from contaminated produce and shellfish
    \item \textbf{Parasites:} giardia and cryptosporidium from contaminated water and food
    \item \textbf{Cross-contamination:} the transfer of organisms from raw meat and poultry to ready-to-eat food via hands, boards, and utensils
    \item \textbf{The danger zone:} bacteria multiply rapidly between 5 and 60 degrees; food left in this zone for more than 2 hours becomes unsafe
    \item \textbf{Cooking temperatures:} thorough cooking kills harmful organisms; poultry must reach 74 degrees in the centre, and minced meat 71 degrees
    \item \textbf{Refrigeration:} keeping the fridge below 5 degrees slows bacterial growth
    \item \textbf{The role of hygiene:} hand washing is the single most effective measure against viral spread
    \item \textbf{Symptoms of illness:} organisms multiply in the gut or produce toxins, causing vomiting, diarrhoea, fever, and cramps
\end{itemize}

\section*{Clinical Features}
Safe food handling has recognisable features.
\begin{itemize}
    \item Hands washed before cooking, after handling raw meat, and after using the toilet
    \item Separate boards and knives for raw meat and ready-to-eat food
    \item Raw meat stored on the bottom shelf of the fridge, below other foods
    \item Food cooked to safe internal temperatures
    \item Leftovers refrigerated within 2 hours and reheated thoroughly
    \item The fridge at or below 5 degrees and the freezer below minus 18 degrees
    \item Food in date, with damaged or swollen cans discarded
    \item Food thawed in the fridge or microwave, never on the bench
\end{itemize}

\section*{Differential Diagnosis}
Food safety decisions address common risks.
\begin{itemize}
    \item The need to wash raw chicken, which spreads bacteria and is not recommended
    \item The difference between use-by and best-before dates
    \item Whether food left out is safe, judged by the 2-hour rule
    \item Whether reheating makes food safe, which it does only if food was chilled promptly
    \item The safety of raw and lightly cooked eggs, which are unsafe for vulnerable groups
    \item The high-risk foods to avoid in pregnancy: deli meats, soft cheeses, raw sprouts, and raw seafood
    \item Food prepared ahead for events, which needs careful temperature control
\end{itemize}

\section*{When to Seek Medical Care}
Food safety is preventive, but foodborne illness needs medical care when it is severe or affects vulnerable people. See a doctor for persistent vomiting, dehydration, blood in the stool, high fever, or illness in a baby, older person, pregnant woman, or immunocompromised person.

\textbf{Call 000 (or local emergency services) for:}
\begin{itemize}
    \item Severe dehydration with confusion, collapse, or no urine
    \item Blood in vomit or stool
    \item Severe abdominal pain with fever
    \item Difficulty breathing
\end{itemize}

\section*{Diagnosis and Investigations}
\begin{itemize}
    \item \textbf{Kitchen review:} checking hand hygiene, storage, cooking, and cooling practices
    \item \textbf{Refrigerator check:} temperature, storage of raw and cooked foods, and use-by dates
    \item \textbf{Cooking check:} using a thermometer to verify safe internal temperatures
    \item \textbf{Food recall awareness:} checking recalls through Food Standards Australia New Zealand for household products
    \item \textbf{Public health advice:} guidelines for safe food handling in pregnancy and for vulnerable groups
\end{itemize}

\section*{Management}
\begin{itemize}
    \item \textbf{Clean:} wash hands, surfaces, and utensils often, and wash fruit and vegetables
    \item \textbf{Separate:} keep raw meat, poultry, and seafood apart from ready-to-eat food
    \item \textbf{Cook:} cook to safe temperatures, using a thermometer, and reheat leftovers until steaming
    \item \textbf{Chill:} refrigerate promptly, keep the fridge below 5 degrees, and never leave perishable food out more than 2 hours
    \item \textbf{Store safely:} raw meat below other foods, covered, and food in sealed containers
    \item \textbf{Thaw safely:} in the fridge, in cold water, or in the microwave, never on the bench
    \item \textbf{Handle leftovers:} refrigerate within 2 hours, eat within a few days, and reheat thoroughly
    \item \textbf{Special care in pregnancy:} avoid high-risk foods and follow listeria-safe advice
\end{itemize}

\section*{Complications}
\begin{itemize}
    \item Foodborne illness from failure of the four rules
    \item Serious infection in vulnerable groups, including listeria in pregnancy
    \item Outbreaks spreading through households and communities
    \item Chronic complications in a minority of infections
    \item Food waste from spoilage, and the cost of unsafe storage
\end{itemize}

\section*{Prognosis and Outcomes}
The rules of food safety prevent nearly all foodborne illness when consistently applied. Households that clean, separate, cook, and chill reduce their risk of food poisoning dramatically, and food prepared safely is safe for everyone, including the most vulnerable. The habits are simple and quickly become automatic. The outcome of good food safety is not just avoiding illness: it is the confidence that the food on the table is safe, every time.

\section*{Prevention and Self-Care}
\begin{itemize}
    \item Wash hands with soap and water before cooking and after handling raw meat
    \item Use separate boards and knives for raw and cooked food
    \item Cook poultry, minced meat, and sausages to a safe internal temperature
    \item Refrigerate leftovers within 2 hours
    \item Keep the fridge at 5 degrees or below
    \item Do not wash raw chicken
    \item Follow the high-risk food advice during pregnancy
    \item Check food recalls and use-by dates, and discard doubtful food
\end{itemize}

\section*{Sources and Further Reading}
The following reputable sources provide further information for healthcare students and patients seeking reliable, current advice about food safety.
\begin{itemize}
    \item Food Safety Information Council. \textit{Food safety tips.} https://foodsafety.asn.au/
    \item Food Standards Australia New Zealand. \textit{Food safety at home.} https://www.foodstandards.gov.au/
    \item Healthdirect Australia. \textit{Food safety.} https://www.healthdirect.gov.au/
    \item NSW Food Authority. \textit{Home food safety.} https://www.foodauthority.nsw.gov.au/
    \item Centers for Disease Control and Prevention. \textit{Four steps to food safety.} https://www.cdc.gov/
    \item World Health Organization. \textit{Five keys to safer food.} https://www.who.int/
\end{itemize}
